\documentclass[b5j,fleqn]{jarticle}
\usepackage{zennyakuUniv}
\begin{document}
\body{
\Title{newcrown3 p57}
\para

\sent{Yesterday, I read the same manga in Japanese and English.}
{昨日、私は同じ漫画を日本語と英語で読みました。}

\sent{I found some pages that have interesting translations of sounds and unique phrases.}
{音や独特な表現が興味深く翻訳されているページをいくつか見つけました。}

\para

\sent{For example, how do you translate a clapping sound?}
{例えば、拍手の音をどのように訳すでしょうか。}

\sent{In Japanese, it's "pachi pachi."}
{日本語では「パチパチ」です。}

\sent{In English, it's "clap clap."}
{英語では「clap clap」です。}

\para

\sent{What about an expression like "Yoroshiku onegaishimasu"?}
{「よろしくお願いします」のような表現はどうでしょうか。}

\sent{There isn't an exact translation for it in English, so it's replaced with phrases like "Good luck."}
{それにぴったり当てはまる英訳はないため、「Good luck」のような表現に置き換えられます。}

\para

\sent{Translations like these help everyone enjoy manga.}
{このような翻訳のおかげで、誰もが漫画を楽しむことができます。}
}
\end{document}
